\chapter{Fazit und Ausblick}\label{fazit-und-ausblick}

Die Diskussion in Kapitel~\ref{diskussion} hat die Evaluationsbefunde eingeordnet und die Grenzen des Systems transparent gemacht. Dieses Kapitel schließt die Arbeit ab. Es umfasst die Beantwortung der Forschungsfrage (Abschnitt~\ref{beantwortung-der-forschungsfrage}), einen strukturierten Ausblick auf kurzfristige, mittelfristige und langfristige Erweiterungen (Abschnitt~\ref{ausblick}) sowie ein abschließendes Schlusswort (Abschnitt~\ref{schlusswort}).



\section{Beantwortung der
Forschungsfrage}\label{beantwortung-der-forschungsfrage}

Die in Abschnitt~\ref{forschungsfrage} formulierte Forschungsfrage lautet:

\begin{quote}
\emph{Lässt sich ein hybrides KI-gestütztes Redaktionssystem für die
Liveticker-Erstellung im Profifußball realisieren, das
Echtzeit-Verarbeitung, Mehrsprachigkeit und White-Label-Betrieb
unterstützt, und inwiefern erfüllt ein solches System journalistische
Qualitätsanforderungen hinsichtlich Korrektheit, Tonalität und
Vollständigkeit?}
\end{quote}

Die Forschungsfrage ist in beiden Teilfragen positiv zu beantworten.

\textbf{Realisierbarkeit:} Das System ist vollständig implementiert und produktionsnah deployt. 23 der 24 definierten Anforderungen sind vollständig erfüllt (vgl.\ Abschnitt~\ref{anforderungsabgleich-soll-ist-vergleich}). F7 (Mehrsprachige Textgenerierung) ist funktional implementiert und durch 21 spanischsprachige Einträge verifiziert. Eine qualitative Evaluation der Übersetzungsqualität steht als Folgestudie aus (vgl.\ Abschnitt~\ref{ausblick}). Die technische Absicherung umfasst 439 automatisierte Tests über drei Schichten sowie eine TypeScript-Typenabdeckung von 98{,}02\,\% der Frontend-Quelldateien (vgl.\ Abschnitt~\ref{technische-qualitätssicherung}). Im Sinne des Design Science Research (vgl. Hevner et al.~2004) leistet die Arbeit einen architektonischen, einen methodischen und einen praktischen Beitrag (vgl. Abschnitt~\ref{beitrag-zur-forschung}).

\textbf{Journalistische Qualität:} Die Textqualität erfüllt die Anforderungen in den Dimensionen Korrektheit (Ø~4{,}6/5) und Vollständigkeit (Ø~4{,}3/5) zuverlässig. Bei der Tonalität ist das Bild differenziert: In der manuellen Evaluation erreicht die Dimension Ø~4{,}1/5 mit klar unterscheidbaren Stilprofilen (vgl.\ Abschnitt~\ref{qualitative-analyse-der-generierten-texte}). Im LLM-as-Judge-Verfahren differenzieren sich die Profile deutlich: \texttt{neutral} erreicht Ø~4{,}76, \texttt{euphorisch} Ø~4{,}11, \texttt{kritisch} Ø~3{,}82. Das kritische Profil zeigt dabei die größte Schwäche bei Tonalität (3{,}50/5) und Vollständigkeit (3{,}62/5). Die höheren Anforderungen des analytisch-bewertenden Stils machen das Prompt-Design zum priorisierten Verbesserungsbedarf (vgl.\ Abschnitt~\ref{ausblick}). Die Prüfaufgabe des Redakteurs verschiebt sich damit von der Faktenkontrolle zur Tonalitätskontrolle, eine Aufgabenteilung, die dem Konzept der „augmented authorship" entspricht (vgl. Bluhm \& Schäfer 2023, S.~35). Ein kontrollierter Feldvergleich bleibt als Folgestudie vorgesehen (vgl. Abschnitt~\ref{ausblick}).

Die modusspezifische Betrachtung ergibt eine klare Empfehlung: Der \texttt{coop}-Modus erweist sich als geeigneter Betriebsmodus für den produktiven Einsatz, da er Beschleunigung und Qualitätssicherung durch Aufgabenteilung verbindet (vgl. Abschnitt~\ref{diskussion-der-betriebsmodi}).



Die identifizierten Qualitätsgrenzen und offenen Validierungsfragen bilden die Grundlage für den folgenden strukturierten Ausblick.

\section{Ausblick}\label{ausblick}

Die Evaluationsbefunde und die Diskussion legen folgende Erweiterungsrichtungen nahe, geordnet nach Umsetzungsnähe.

\paragraph{Kurzfristig}
\begin{itemize}\tightlist
\item \textbf{Authentifizierung und Rollensystem:} Die Integration eines JSON Web Token (JWT)- oder OAuth-2.0-basierten (Open Authorization) Authentifizierungsframeworks ist die wichtigste Voraussetzung für einen Mehrbenutzerbetrieb. Ein rollenbasiertes Zugriffskonzept (Redakteur, Chefredakteur, Administrator) würde die Freigabe-Workflows im \texttt{coop}-Modus um eine Vier-Augen-Prüfung erweitern.
\item \textbf{Prompt-Optimierung für das kritische Stilprofil:} Das kritische Stilprofil erzielt im LLM-as-Judge-Verfahren die niedrigsten Werte bei Tonalität (3{,}50/5) und Vollständigkeit (3{,}62/5, vgl.\ Abschnitt~\ref{erweiterte-evaluation}). Die höheren Anforderungen des analytisch-bewertenden Stils werden durch die bestehenden Stilreferenzen, die überwiegend euphorisch ausgerichtet sind, unzureichend abgedeckt (vgl.\ Abschnitt~\ref{stilreferenzen}). Gezielte Anpassungen mit stilgefilterten kritischen Referenzen und präzisierten Stilinstruktionen würden diesen Befund direkt adressieren.
\item \textbf{Diversitätsmechanismen gegen Phrasenwiederholung:} 13\,\% der manuell bewerteten Einträge zeigten formelhaften Satzbau bei aufeinanderfolgenden Ereignissen gleichen Typs (vgl.\ Abschnitt~\ref{qualitative-analyse-der-generierten-texte}). Techniken wie Negative Prompting (explizites Verbot zuletzt verwendeter Formulierungen) oder eine ereignistypspezifische Varianten-Bibliothek würden formelhaften Satzbau zuverlässig unterbinden.
\item \textbf{Stilspezifische Temperaturkonfiguration:} Der Temperaturvergleich (vgl. Abschnitt~\ref{temperatur-context}) zeigt, dass $T\,=\,0{,}3$ für das \texttt{euphorisch}-Profil und $T\,=\,0{,}7$ für \texttt{neutral} und \texttt{kritisch} die jeweils besten Qualitätswerte liefern. Eine stilspezifische Temperaturkonfiguration ist ohne Architekturänderungen durch einfache Parametrierung umsetzbar.
\item \textbf{Streaming-Optimierung:} Die gemessene LLM-Latenz (Median 3\,338~ms, vgl. Abschnitt~\ref{llm-latenz}) ließe sich durch Token-Streaming reduzieren: Generierte Texte würden dem Redakteur inkrementell angezeigt, statt erst nach vollständiger Generierung.
\item \textbf{n8n-Workflow-Tests:} Die 17 n8n-Workflows sind derzeit ausschließlich manuell verifiziert (vgl.\ Abschnitt~\ref{kritische-reflexion}). Integrationstests gegen eine Staging-Datenbank würden die Zuverlässigkeit der Datenversorgungspipeline systematisch absichern.
\end{itemize}

\paragraph{Mittelfristig}
\begin{itemize}\tightlist
\item \textbf{Kontrollierter Feldvergleich mit externer Qualitätsevaluation:} Eine Folgestudie mit professionellen Sportredakteuren unter Echtbedingungen würde zwei bisher offene Fragen schließen. Erstens ließen sich empirisch belastbare Aussagen zu Time-to-Publish, Fehlerrate und Redakteurszufriedenheit je Betriebsmodus (\texttt{manual}, \texttt{coop}, \texttt{auto}) gewinnen. Zweitens würde sie eine externe Validierung der Textqualität durch unabhängige Rater ermöglichen. Die externen Rater bewerten bestehende KI-generierte Einträge blind und lösen damit die Selbstevaluationsproblematik der Hauptevaluation auf. Die implementierte Evaluationsinfrastruktur (TTP-Messung, Cohen's~Kappa, Cliff's~Delta) ist für beide Studienteile bereits vorbereitet.
\item \textbf{Evaluation der Mehrsprachigkeit:} Eine separate Qualitätsevaluation der mehrsprachigen Textgenerierung, insbesondere für die im Kontext von Eintracht Frankfurt relevante Sprache Englisch, würde die in Abschnitt~\ref{mehrsprachigkeit-als-herausforderung} motivierte Internationalisierungsanforderung empirisch absichern.
\item \textbf{Fine-Tuning statt Few-Shot:} Für vereinsspezifische Instanzen könnte ein Fine-Tuning auf dem Korpus historischer Vereinsticker die stilistische Konsistenz über den Few-Shot-Ansatz hinaus verbessern. Die \texttt{style\_references}-Tabelle bildet bereits eine kuratierte Datenbasis, die als Ausgangspunkt für ein Fine-Tuning-Dataset dienen könnte.
\end{itemize}

\paragraph{Langfristig}
\begin{itemize}\tightlist
\item \textbf{Multimodale Ticker:} Eine Folgestudie könnte klären, ob Vision-Language-Modelle aus Spielszenenbildern (z.\,B. über die bestehende ScorePlay-Integration) ereignisrelevante Bildbeschreibungen generieren, die mit den strukturierten API-Ereignisdaten kohärent kombinierbar sind. Ein kontrollierter Vergleich zwischen faktenbasierter Textgenerierung (aktueller Ansatz) und VLM-basierter Bildauswertung würde erstmals die Qualitätsdifferenz beider Ansätze anhand der drei Qualitätsdimensionen dieser Arbeit messen.
\item \textbf{Autonomes Scoring der Textqualität:} Die Ablösung der manuellen Qualitätsbewertung durch ein trainiertes Scoring-Modell, das Korrektheit, Tonalität und Vollständigkeit automatisch bewertet, würde eine kontinuierliche Qualitätsüberwachung im Produktivbetrieb ermöglichen. Die implementierte \texttt{aggregate\_quality\_by\_group}-Funktion deckt die erforderliche Aggregationslogik bereits ab.
\item \textbf{Generalisierung über den Fußball hinaus:} Die Provider-agnostische LLM-Architektur und das Event-basierte Generierungsmodell sind nicht fußballspezifisch. Eine Erweiterung auf andere Sportarten (Handball, Basketball, Eishockey) erforderte primär die Anpassung der Event-Typ-Mappings, der Phasendefinitionen und der Kontextbuilder. Die Kernarchitektur würde unverändert bleiben.
\end{itemize}



\section{Schlusswort}\label{schlusswort}

Die vorliegende Arbeit zeigt, dass ein hybrides KI-gestütztes
Redaktionssystem den operativen Zeitdruck bei der Liveticker-Erstellung
bewältigen kann, ohne die redaktionelle Kontrolle aufzugeben. Der
\texttt{coop}-Modus realisiert das in Abschnitt~\ref{problemstellung-und-zielsetzung} formulierte Zielbild
eines Systems, in dem „die finale Entscheidungshoheit und publizistische
Verantwortung beim Menschen verbleiben". Die White-Label-Architektur adressiert den Bedarf an vereinsspezifischer Identität ohne Systemduplizierung. Die mehrsprachige Generierungskomponente legt die technische Grundlage für die Internationalisierung, deren Qualitätssicherung als nächster Schritt aussteht.

Die KI-gestützte Liveticker-Generierung ersetzt den Redakteur nicht. Sie reduziert den kognitiven Engpass der simultanen Textproduktion unter Echtzeit-Bedingungen und schafft damit Kapazität für journalistische Kernleistungen: kontextuelle Einordnung, kritische Bewertung und das Wahrnehmen redaktioneller Verantwortung. Die von Constantin Seibt beschriebene „komplette Überforderung" des Liveticker-Autors (Seibt, zit.~n. Beils 2023, S.~56) bleibt eine strukturelle Bedingung des Genres. Das System unterstützt den Autor dabei, ohne die Entscheidung abzunehmen.
